\paragraph{平方根谱碰撞与前导零点的 \texorpdfstring{$n^{-2}$}{n^{-2}} 量子化}
令 \(U\subset\CC\) 为开集，\(A:U\to\Mat_{d}(\CC)\) 为解析矩阵族（\(d\ge 2\)）。
取非零线性泛函 \(\ell\in(\CC^d)^\vee\) 与非零向量 \(r\in\CC^d\)，并定义有限尺度配分函数
\[
Z_n(u):=\ell\!\bigl(A(u)^n r\bigr),\qquad n\in\NN.
\]
记特征多项式
\[
P(u,\lambda):=\det(\lambda I-A(u)).
\]

\begin{theorem}[平方根谱碰撞下 Lee--Yang 前导零点的 \texorpdfstring{$n^{-2}$}{n^{-2}} 量子化定标]\label{thm:xi-leyang-square-root-collision-leading-zeros-n2}
\leanverified{paper\_xi\_leyang\_square\_root\_collision\_leading\_zeros\_n2}
设存在 \((u_c,\lambda_c)\in U\times\CC\) 满足：
\begin{enumerate}
  \item \(P(u_c,\lambda_c)=0\) 且 \(\partial_\lambda P(u_c,\lambda_c)=0\)；
  \item \(\partial_{\lambda\lambda}P(u_c,\lambda_c)\neq 0\) 且 \(\partial_u P(u_c,\lambda_c)\neq 0\)；
  \item \(A(u_c)\) 的其余全部特征值 \(\{\lambda_j(u_c)\}_{j\ge 3}\) 满足 \(|\lambda_j(u_c)|<|\lambda_c|\)；
  \item \(\ell,r\) 在该二重特征值对应的二维广义谱簇上诱导的两支谱投影系数不同时为零。
\end{enumerate}
则存在 \(\varepsilon>0\) 及解析函数 \(c_\pm(t)\)（满足 \(c_\pm(0)\neq 0\)），以及定义于 \(|t|<\varepsilon\) 的两支特征值分支 \(\lambda_\pm(t)\)，使得在二重覆盖参数化
\[
u=u_c+t^2
\]
下成立展开
\[
\lambda_\pm(t)=\lambda_c\pm \alpha\,t+O(t^2),
\qquad
\alpha^2=-\frac{2\,\partial_u P(u_c,\lambda_c)}{\partial_{\lambda\lambda}P(u_c,\lambda_c)}\neq 0,
\]
并且存在常数 \(\eta\in(0,|\lambda_c|)\) 使对 \(|t|<\varepsilon\) 一致有
\[
Z_n(u_c+t^2)=c_+(t)\lambda_+(t)^n+c_-(t)\lambda_-(t)^n+O(\eta^n)
\qquad(n\to\infty).
\]
令
\[
\kappa:=-\frac{c_-(0)}{c_+(0)}\in\CC^\times,
\]
并固定一条对数分支 \(\log\) 在某邻域内解析。
则存在常数 \(C>0\) 使得对每个固定整数 \(k\in\ZZ\)，当 \(n\) 足够大时，方程 \(Z_n(u)=0\) 在缩放窗 \(|u-u_c|\le Cn^{-2}\) 内存在唯一零点 \(u_{n,k}\)，且满足
\[
u_{n,k}
=
u_c+\frac{\lambda_c^{2}}{4\alpha^{2}n^{2}}
\bigl(\log\kappa+2\pi i k\bigr)^{2}
+O(n^{-3}),
\qquad(n\to\infty).
\]
\end{theorem}
\begin{proof}
由 (1)(2) 与 Weierstrass 制备定理，\(P(u,\lambda)\) 在 \((u_c,\lambda_c)\) 的邻域内可分解为
\[
P(u,\lambda)=\Bigl((\lambda-\lambda_c)^2+a_1(u)(\lambda-\lambda_c)+a_0(u)\Bigr)\,Q(u,\lambda),
\qquad Q(u_c,\lambda_c)\neq 0,
\]
其中 \(a_0(u_c)=a_1(u_c)=0\)，且 \(a_0'(u_c)\neq 0\)（横截性）。
因此存在二重覆盖 \(u=u_c+t^2\)，使得该二次因子的两根成为解析分支 \(\lambda_\pm(t)\)，并由二次项比较得到
\[
(\lambda-\lambda_c)^2
=
-\frac{2\,\partial_u P(u_c,\lambda_c)}{\partial_{\lambda\lambda}P(u_c,\lambda_c)}(u-u_c)
+O\bigl(|u-u_c|^{3/2}\bigr),
\]
从而 \(\lambda_\pm(t)=\lambda_c\pm\alpha t+O(t^2)\) 且 \(\alpha^2\) 如述。

由解析谱投影与 (3) 的严格谱隙，可将 \(A(u_c+t^2)^n\) 分解为两主支谱投影项与余项：
\[
A(u_c+t^2)^n=\lambda_+(t)^n\Pi_+(t)+\lambda_-(t)^n\Pi_-(t)+R_n(t),
\]
其中 \(\Pi_\pm(t)\) 解析，且 \(\|R_n(t)\|=O(\eta^n)\) 对 \(|t|<\varepsilon\) 一致成立。
左右夹以 \(\ell(\cdot r)\) 得
\(
Z_n(u_c+t^2)=c_+(t)\lambda_+(t)^n+c_-(t)\lambda_-(t)^n+O(\eta^n)
\)，其中 \(c_\pm(t):=\ell(\Pi_\pm(t)r)\)；
由 (4) 得 \(c_\pm(0)\neq 0\)。

令 \(Z_n(u_c+t^2)=0\)。
将余项并入误差后，有
\[
\Bigl(\frac{\lambda_+(t)}{\lambda_-(t)}\Bigr)^n=-\frac{c_-(t)}{c_+(t)}+O\bigl((\eta/|\lambda_c|)^n\bigr).
\]
在 \(|t|<\varepsilon\) 内取对数并吸收分支不定性 \(2\pi i k\) 得
\[
n\log\frac{\lambda_+(t)}{\lambda_-(t)}
=
\log\kappa+2\pi i k+O(t)+O\bigl((\eta/|\lambda_c|)^n\bigr).
\]
由 \(\lambda_\pm(t)=\lambda_c\pm\alpha t+O(t^2)\) 得
\(
\log\frac{\lambda_+(t)}{\lambda_-(t)}=\frac{2\alpha}{\lambda_c}t+O(t^2)
\)。
因此
\[
t=\frac{\lambda_c}{2\alpha n}\bigl(\log\kappa+2\pi i k\bigr)+O(n^{-2}),
\]
进而
\(
u-u_c=t^2=\frac{\lambda_c^2}{4\alpha^2n^2}(\log\kappa+2\pi i k)^2+O(n^{-3})
\)。
唯一性可由上述对数方程在 \(|t|\le c/n\) 区域的收缩估计得到。证毕。
\end{proof}

\begin{corollary}[对称可见性下的两零点外推与幅度消去]\label{cor:xi-leyang-two-leading-zeros-extrapolate-uc}
在定理 \ref{thm:xi-leyang-square-root-collision-leading-zeros-n2} 的假设下，若进一步有 \(c_+(0)=c_-(0)\)，则 \(\kappa=-1\)。
取对数分支满足 \(\log(-1)=i\pi\)，并令
\[
C:=\frac{\pi^2\lambda_c^2}{4\alpha^2}\neq 0.
\]
则当 \(n\to\infty\) 时，靠近 \(u_c\) 的两枚最小尺度零点（对应 \(k=0,1\)）满足
\[
u_{n,0}=u_c-\frac{C}{n^2}+O(n^{-3}),
\qquad
u_{n,1}=u_c-\frac{9C}{n^2}+O(n^{-3}).
\]
从而两零点外推
\[
\widehat u_c(n):=\frac{9u_{n,0}-u_{n,1}}{8}
=u_c+O(n^{-3}),
\]
并且幅度可由
\[
\widehat C(n):=\frac{n^2}{8}\bigl(u_{n,0}-u_{n,1}\bigr)=C+O(n^{-1})
\]
一致估计。
\par
同样地，幅度亦可由两零点外推出的消元式直接读出：
\[
\widehat C_{\mathrm{elim}}(n):=-n^2\bigl(u_{n,0}-\widehat u_c(n)\bigr)=C+O(n^{-1}).
\]
\leanverified{paper\_xi\_leyang\_two\_leading\_zeros\_extrapolate\_uc}
\end{corollary}
\begin{proof}
由定理 \ref{thm:xi-leyang-square-root-collision-leading-zeros-n2} 的量子化公式代入 \(\log\kappa=\log(-1)=i\pi\)，得到
\[
u_{n,k}
=u_c+\frac{\lambda_c^2}{4\alpha^2n^2}\bigl(i\pi+2\pi i k\bigr)^2+O(n^{-3})
=u_c-\frac{\pi^2\lambda_c^2}{4\alpha^2n^2}(2k+1)^2+O(n^{-3}).
\]
取 \(k=0,1\) 即得前两式；线性消元得到 \(\widehat u_c(n)\) 的表达式，\(\widehat C(n)\) 由差分读出。证毕。
\end{proof}

\begin{theorem}[三零点端点外推的 \texorpdfstring{$O(n^{-6})$}{O(n^-6)} 误差改进]\label{thm:xi-leyang-three-leading-zeros-extrapolate-uc-n6}
\leanverified{paper\_xi\_leyang\_three\_leading\_zeros\_extrapolate\_uc\_n6}
设在某端点机制下，存在常数 \(u_c\in\CC\)、\(C\neq 0\)、\(D\in\CC\)，使得对固定的 \(k\in\{0,1,2\}\) 有渐近展开
\[
u_{n,k}
=
u_c
-\frac{C}{n^{2}}(2k+1)^{2}
-\frac{D}{n^{4}}(2k+1)^{4}
O(n^{-6})
\qquad(n\to\infty).
\]
定义三零点外推量
\[
\boxed{\
\widehat u_{c}^{(n)}
:=\frac{75}{64}\,u_{n,0}-\frac{25}{128}\,u_{n,1}+\frac{3}{128}\,u_{n,2}.
}
\]
则成立严格误差改进
\[
\boxed{\ \widehat u_{c}^{(n)}=u_{c}+O(n^{-6}).\ }
\]
\end{theorem}
\begin{proof}
令 \(a_k:=(2k+1)^2\)，则 \(a_0=1,a_1=9,a_2=25\)。
取系数 \(\alpha_0,\alpha_1,\alpha_2\) 满足
\[
\alpha_0+\alpha_1+\alpha_2=1,\qquad
\alpha_0a_0+\alpha_1a_1+\alpha_2a_2=0,\qquad
\alpha_0a_0^2+\alpha_1a_1^2+\alpha_2a_2^2=0.
\]
该 Vandermonde 线性系统唯一解为
\(
(\alpha_0,\alpha_1,\alpha_2)=(75/64,\,-25/128,\,3/128)
\)。
将展开式代入 \(\sum_{k=0}^2\alpha_k u_{n,k}\)，由约束恰消去 \(n^{-2}\) 与 \(n^{-4}\) 两阶项，余项阶为 \(O(n^{-6})\)。证毕。
\end{proof}

\begin{theorem}[\texorpdfstring{$(M+1)$}{M+1} 点端点外推的闭式权重与最优误差阶]\label{thm:xi-leyang-mplus1-point-extrapolate-uc-optimal}
\leanverified{paper\_xi\_leyang\_mplus1\_point\_extrapolate\_uc\_optimal}
在与定理 \ref{thm:xi-leyang-three-leading-zeros-extrapolate-uc-n6} 同型的端点机制下，固定整数 \(M\ge 1\)，并令
\[
a_k:=(2k+1)^2\qquad (k=0,1,\dots,M).
\]
定义 Lagrange 权重
\[
\boxed{\
\alpha_{k}^{(M)}
:=
\prod_{\substack{0\le i\le M\\ i\ne k}}
\frac{-a_{i}}{a_{k}-a_{i}}
\qquad(k=0,\dots,M),
}
\]
以及外推估计量
\[
\boxed{\
\widehat u_{c,M}^{(n)}:=\sum_{k=0}^{M}\alpha_{k}^{(M)}\,u_{n,k}.
}
\]
若存在常数 \(C_1,\dots,C_{M+1}\in\CC\) 使得对每个固定 \(k\in\{0,1,\dots,M\}\) 有
\[
u_{n,k}
=
u_c+\sum_{j=1}^{M+1}\frac{C_j}{n^{2j}}\,a_k^{j}+O(n^{-2M-4}),
\qquad(n\to\infty),
\]
则
\[
\boxed{\ \widehat u_{c,M}^{(n)}=u_c+O(n^{-2M-2}).\ }
\]
并且在“仅使用 \(\{u_{n,0},\dots,u_{n,M}\}\) 并要求消去前 \(M\) 个偶阶项”的信息模型下，上述权重是唯一解，从而达到该模型内的最优误差阶。
\end{theorem}
\begin{proof}
在假设展开中，截断到 \(n^{-2M-2}\) 的部分在变量 \(a_k\) 上是次数至多 \(M+1\) 的多项式，并且 \(u_c\) 对应于 \(a=0\) 处的取值。
因此 \(\widehat u_{c,M}^{(n)}\) 等价于对节点 \(\{a_0,\dots,a_M\}\) 的 Lagrange 插值多项式在 \(a=0\) 的取值，其权重即为 \(L_k(0)\)，从而给出闭式 \(\alpha_k^{(M)}\)。
插值误差由最高项 \(a^{M+1}\) 控制，对应尺度 \(n^{-2M-2}\)，得到误差阶。
唯一性来自 Vandermonde 矩阵非退化。证毕。
\end{proof}

\begin{corollary}[由三零点同时恢复主标度常数与次阶常数]\label{cor:xi-leyang-recover-C-D-from-three-leading-zeros}
\leanverified{paper\_xi\_leyang\_recover\_C\_D\_from\_three\_leading\_zeros}
在定理 \ref{thm:xi-leyang-three-leading-zeros-extrapolate-uc-n6} 的假设下，令 \(\widehat u_c^{(n)}\) 为其三零点外推量，并定义
\[
B_k:=n^{2}\bigl(\widehat u_{c}^{(n)}-u_{n,k}\bigr)\qquad(k=0,1).
\]
则常数 \(C,D\) 可由显式公式一致估计：
\[
\boxed{\
\widehat C^{(n)}:=\frac{81B_{0}-B_{1}}{72}
\ =\ C+O(n^{-4}),\ }
\qquad
\boxed{\
\widehat D^{(n)}:=\frac{n^{2}(B_{1}-9B_{0})}{72}
\ =\ D+O(n^{-2}).\ }
\]
\end{corollary}
\begin{proof}
由定理 \ref{thm:xi-leyang-three-leading-zeros-extrapolate-uc-n6}，
\(
\widehat u_c^{(n)}=u_c+O(n^{-6})
\)。
结合 \(a_0=1,a_1=9\) 与展开式可得
\[
B_k=C\,a_k+\frac{D}{n^2}\,a_k^2+O(n^{-4})\qquad(k=0,1).
\]
对 \(k=0,1\) 作线性消元：\(\widehat C^{(n)}\) 消去 \(D/n^2\) 项而达到 \(O(n^{-4})\)，\(\widehat D^{(n)}\) 由差分提取 \(D\) 并得到 \(O(n^{-2})\)。证毕。
\end{proof}

\begin{proposition}[平方根分歧点的等模曲线切线方向]\label{prop:xi-leyang-equimodular-tangent-square-root}
在定理 \ref{thm:xi-leyang-square-root-collision-leading-zeros-n2} 的记号下，定义局部等模集合
\[
\Sigma:=\{u\in U:\ |\lambda_+(u)|=|\lambda_-(u)|\},
\]
其中 \(\lambda_\pm\) 取定理中的两支 Puiseux 解析延拓。
则在足够小邻域内，\(\Sigma\) 由恰好两条穿过 \(u_c\) 的实解析弧组成；
其在 \(u\)-平面中穿过 \(u_c\) 的切线方向唯一确定为
\[
\arg(u-u_c)\equiv \pi-2\arg\bigl(\lambda_c\overline{\alpha}\bigr)\pmod{\pi}.
\]
更精确地，在均匀化坐标 \(u=u_c+t^2\) 下，等模条件的一阶近似等价于
\[
\Re\bigl(\lambda_c\overline{\alpha}\,t\bigr)=0,
\]
从而在 \(t\)-平面给出两条互逆直线，其经 \(t\mapsto t^2\) 推送为上述两条切线方向。
\end{proposition}
\begin{proof}
由 \(\lambda_\pm(t)=\lambda_c\pm\alpha t+O(t^2)\) 得
\[
|\lambda_+(t)|^2-|\lambda_-(t)|^2
=|\lambda_c+\alpha t|^2-|\lambda_c-\alpha t|^2+O(|t|^2)
=4\,\Re(\lambda_c\overline{\alpha}\,t)+O(|t|^2).
\]
因此 \(|\lambda_+|=|\lambda_-|\) 在一阶上等价于 \(\Re(\lambda_c\overline{\alpha}\,t)=0\)，给出 \(t\)-平面中过原点的两条直线。
由 \(u-u_c=t^2\) 切向角加倍，得到所示同余式；高阶项由实解析隐函数定理给出两条实解析弧。证毕。
\end{proof}

\begin{theorem}[多项式传递矩阵的临界点代数性与判别式复杂度上界]\label{thm:xi-polynomial-transfer-matrix-branchpoint-discriminant-degree}
设 \(A(u)\in\Mat_d(\ZZ[u])\) 为多项式矩阵，且每个矩阵元 \(A_{ij}(u)\) 的次数不超过 \(k\ge 1\)。
令
\[
P(u,\lambda):=\det(\lambda I-A(u))\in\ZZ[u,\lambda],
\qquad
D(u):=\mathrm{Disc}_\lambda P(u,\lambda)\in\ZZ[u].
\]
则成立：
\begin{enumerate}
  \item \(D(u)\) 非恒零，且
  \[
  \deg_u D\le (2d-1)\deg_u P\le (2d-1)dk.
  \]
  \item 任一谱分歧点 \(u_c\)（即存在 \(\lambda\) 使 \(P(u_c,\lambda)=\partial_\lambda P(u_c,\lambda)=0\)）必为代数数，且
  \[
  [\QQ(u_c):\QQ]\le \deg_u D \le (2d-1)dk.
  \]
  \item 若某分歧点 \(u_c\) 为 \(D\) 的单根（\(D(u_c)=0\) 且 \(D'(u_c)\neq 0\)），并且在该点存在严格谱隙 \(|\lambda_j(u_c)|<|\lambda_c|\)（除去发生碰撞的二重根 \(\lambda_c\)），则定理 \ref{thm:xi-leyang-square-root-collision-leading-zeros-n2} 的平方根碰撞情形在该点为代数可审计：可由判别式单根性与谱隙验证推出 \(n^{-2}\) 量子化定标。
\end{enumerate}
\end{theorem}
\begin{proof}
由 \(D(u)=\mathrm{Disc}_\lambda(P)=(-1)^{d(d-1)/2}\mathrm{Res}_\lambda(P,\partial_\lambda P)\)。
对任意 \(F,G\in\ZZ[u,\lambda]\)，其关于 \(\lambda\) 的结式满足次数估计
\[
\deg_u\mathrm{Res}_\lambda(F,G)
\le (\deg_\lambda F)(\deg_u G)+(\deg_\lambda G)(\deg_u F).
\]
取 \(F=P\)、\(G=\partial_\lambda P\)，则 \(\deg_\lambda P=d\)、\(\deg_\lambda(\partial_\lambda P)=d-1\)，且 \(\deg_u(\partial_\lambda P)\le \deg_u P\)，从而
\[
\deg_u D\le d\,\deg_u P+(d-1)\deg_u P=(2d-1)\deg_u P.
\]
又 \(P(u,\lambda)=\det(\lambda I-A(u))\) 的每一项为 \(d\) 个矩阵元之积，故 \(\deg_u P\le dk\)，得第 1 条。

第 2 条由分歧点满足 \(D(u_c)=0\) 得出：\(u_c\) 为整系数多项式 \(D\) 的零点，因而为代数数且次数不超过 \(\deg_u D\)。

对第 3 条，\(D'(u_c)\neq 0\) 等价于 \(P(u,\lambda)\) 在 \((u_c,\lambda_c)\) 发生单纯二重根碰撞（局部仅两根以平方根方式分歧），与定理 \ref{thm:xi-leyang-square-root-collision-leading-zeros-n2} 的 (1)(2) 同型；
严格谱隙与可见性条件给出两主支主导，从而可应用定理 \ref{thm:xi-leyang-square-root-collision-leading-zeros-n2} 得到 \(n^{-2}\) 量子化。证毕。
\end{proof}

\begin{conjecture}[平方根端点下的两零点比值刚性与超收敛外推]\label{conj:xi-leyang-square-root-endpoint-two-zero-ratio}
设一族有限尺度配分函数 \(Z_n(u)\) 来自有限维解析传递算子或传递矩阵（例如 \(Z_n(u)=\ell(A(u)^n r)\)），并且其自由能的主要解析延拓障碍由某临界端点 \(u_c\) 处两条主导特征值的单纯平方根碰撞控制（判别式在 \(u_c\) 为单根，且存在严格谱隙与可见性）。
则存在编号约定使得靠近 \(u_c\) 的前两枚零点满足
\[
u_{n,1}-u_c=9\,(u_{n,0}-u_c)+o(n^{-2}),
\]
并且由两零点组合
\(
\widehat u_c(n)=(9u_{n,0}-u_{n,1})/8
\)
给出的端点外推误差为 \(o(n^{-2})\)。
若该刚性在数值数据中成立，则可据此反推临界端点处主导谱碰撞为单纯平方根类型，并据判别式与谱隙对端点作代数定位。
\end{conjecture}
